\section{Quantum Slicing de Finetti Theorem}\label{sec:de_finetti}

In this section, we prove \cref{thm:de-Finnetti-type-1}. In spirit, it is a de Finetti type theorem with the contraint that there is little entanglement across some cut. We refer to such type of theorem as the slicing de Finetti theorem.


\subsection{One-versus-Many Slicing de Finetti}
To start, we study the following most basic scenario that a given permutation-invariant pure quantum state from $\Symm{k}{d}$ has a large top Schmidt coefficient over cut between the first and the remaining subsystems. We obtain a dimension independent quantum de Finetti theorem under slicing constraints from first principles.


\begin{theorem}[One-versus-many Slicing de Finetti]\label{thm:one_vs_many}
Let $\ket{\sigma}^{A_1\ldots A_k} \in \vee^k(\C^d)$. If the largest Schmidt coefficient across
the cut $A_1:A_2\cdots A_k$ is at least $\sqrt{1-\epsilon}$, then
  \begin{align*}
    \max_{\ket{\phi} \in \C^d} \abs{\bra{\sigma}^{A_1\ldots A_k} \ket{\phi}^{\otimes k} }^2  \ge 1 - 8k^3\cdot \epsilon.
  \end{align*}
\end{theorem}



To prove this theorem, we first establish the following duplicate lemma. It says that when a symmetric state $\ket\sigma$ is close to some product state $\ket\phi\ket\rho$, then you can find a new state close to $\ket\sigma$ that with two $\ket\phi$ and harms the closeness only mildly.
\begin{lemma}[Duplicate Lemma]\label{lem:dup}
Let $\ket{\sigma}\in \Symm{k}{d}$. 
Consider some arbitrary decomposition of $\{A_1,A_2,\ldots, A_k\}=\cA\cup \cB \cup \cC$,
such that $|\cA|=|\cB|$. Suppose $ |\bra\sigma ^{\cA\cB\cC} \ket\phi^{\cA} \ket\rho^{\cB\cC}|^2 \ge 1-\epsilon. $
Then, there is a state $\ket\zeta^{\cA\cB\cC}$ such that $\ket\zeta=\ket\phi^\cA\ket\phi^\cB$ $\ket\gamma^{C}$ for some $\ket\gamma^C$, and
    \[
        |\langle\sigma\mid \zeta\rangle|^2 \ge 1-8\epsilon.
    \]
Furthermore, if $\rho_\cC$ is a pure state, then $\gamma = \rho_\cC$.
\end{lemma}

\begin{proof}
We assume that $\epsilon<1/8$, otherwise the statement is trivially true.
Apply Schmidt decomposition to $\ket\rho^{\cB\cC}$ for the $\cB:\cC$ cut,
\[
    \ket\rho^{\cB\cC} = \sum_i \sqrt{\lambda_i} \ket{\beta_i}^\cB\ket{\gamma_i}^\cC.
\]
Let 
\[
    \ket{\rho'}^{\cA\cC}=\sum_i \sqrt{\lambda_i} \ket{\beta_i}^\cA\ket{\gamma_i}^\cC.
\]
Since $\ket\sigma\in\Symm{k}{d}$, we have
\[
    |
        \bra\sigma^{\cA\cB\cC} \ket\phi^\cA\ket\rho^{\cB\cC}|^2
    =
    |
        \bra\sigma^{\cA\cB\cC} \ket\phi^\cB\ket{\rho'}^{\cA\cC}
    |^2
    = 1- \epsilon.
\]
By \cref{fact:overlap-triangle-ineq},
\begin{align}
    (1-2\epsilon)^2
        &\le |\bra\phi^\cA\bra\rho^{\cB\cC} \ket\phi^\cB\ket{\rho'}^{\cA\cC}|^2
        = \left( \sum_i \lambda_i |\langle \phi\mid \beta_i \rangle|^2 \right)^2. \label{eq:slice-ineq}
\end{align}
Abbreviate $\eta_i = |\langle\phi\mid \beta_i \rangle|^2$. Note that
\[
    \sum \eta_i  \le 1, \quad \sum \lambda_i = 1.
\]
Therefore, immediately from (\ref{eq:slice-ineq}),  
\begin{align}
    \lambda_1, \max \eta_i \ge 1-2\epsilon,
\end{align}
which is at least 3/4 since $\epsilon<1/8$. If $\eta_1\not=\max\eta_i$, then 
\[
1-2\epsilon \le \sum_i \lambda_i \eta_i \le \lambda_1 (1-\max \eta_i) + \max \eta_i \cdot (1-\lambda_1) \le 4\epsilon,
\]
which is impossible as $\epsilon < 1/8$. Therefore, $\eta_1 \ge 1-2\epsilon$.

We push it further,
\begin{align}
    1-2\epsilon 
    &\le \sum_i \lambda_i \eta_i \le \lambda_1 \eta_1 + (1-\lambda_1)(1-\eta_1) = 2\lambda_1 \eta_1 -\lambda_1-\eta_1 +1
    \nonumber \\
    &\le 2\lambda_1\eta_1 -2\sqrt{\lambda_1\eta_1} +1
    \nonumber \\
    &=2 \left(\sqrt{\lambda_1 \eta_1 }-\frac{1}{2}\right)^2 + \frac{1}{2},
    \nonumber
    %\label{eq:top-top-bound}
\end{align}
where the second step is due to AM-GM inequality. Since $\lambda_1, \eta_1 > 3/4$, and $\epsilon<1/8$, 
\begin{align}
    \lambda_1\eta_1 \ge \left(\frac{1}{2}+\sqrt{\frac{1}{4}-\epsilon}\right)^2 \ge  1-3\epsilon,\nonumber
\end{align}
where the last inequality holds for $\epsilon\in[0,1/8]$. 
% see https://www.wolframalpha.com/input?i=1%2F2%2Bsqrt%281-4x%29%2F2+-+sqrt%281-3x%29
Note that
\[
 |\bra\phi^\cA\bra\rho^{\cB\cC} \ket\phi^\cA\ket\phi^\cB\ket{\gamma_1}^\cC|^2\ge\lambda_1\eta_1 \ge 1-3\epsilon.
\]
By \cref{fact:overlap-triangle-ineq} and that $1-3\epsilon > 1/2$, it can be verified that
\begin{align*}
    |\bra\sigma^{\cA\cB\cC} \ket\phi^\cA\ket\phi^\cB\ket{\gamma_1}^\cC|^2 
    &\ge (\sqrt{(1-\epsilon)(1-3\epsilon)}-\sqrt{3}\epsilon)^2
    \\
    &= 1-4\epsilon+6\epsilon^2 -2\sqrt{3}\epsilon\sqrt{(1-\epsilon)(1-3\epsilon)}
    \\
    &\ge 1-8\epsilon.\qedhere
\end{align*}
\end{proof}

Now \cref{thm:one_vs_many} is a simple consequence of \cref{lem:dup}: Duplicate the the first subsystem taken from the top Schmidt component of $\ket\sigma$.
\begin{proof}[Proof of \emph{\cref{thm:one_vs_many}}]
Let $\ket{\sigma_0}=\ket\phi\ket\gamma$ be the top Schmidt component of $\ket\sigma$ for the $A_1: A_2\ldots A_k$ cut. By assumption of the theorem statement,
\[
    |\langle \sigma \mid \sigma_0\rangle|^2 \ge 1-\epsilon.
\]
Let $m =\lfloor \log k \rfloor, m^* = \lceil \log k \rceil$. For $i = 1, 2, \ldots, m$, apply the Duplicate Lemma on $\ket{\sigma_{i-1}}$ with $\cA = \{A_1,A_2,\ldots, A_{2^{i-1}}\}, \cB=\{A_{2^{i-1}+1}, A_{2^{i-1}+2},\ldots, A_{2^i}\}$. Let $\ket{\sigma_i}$ be the $\ket\zeta$ guaranteed by the Duplicate Lemma. \iffalse
It follows that
\begin{enumerate}
    \item  $|\langle \sigma \mid \sigma_{m}\rangle|^2 \ge 1 - 8^{m}\epsilon$,%\ge 1-8k^3\epsilon$
    \item  $\ket{\sigma_{m}}=\ket\phi^{\otimes k}$.
\end{enumerate}
\fi
If $2^m < k$, apply the Duplicate Lemma one more time on $\ket{\sigma_m}$ with $\cA=\{A_1, A_2,\ldots, A_{k - 2^m}\}, \cB=\{A_{2^m+1}, A_{2^m+2}, \ldots, A_k\}$, and let $\ket{\sigma_{m^*}}$ be the state guaranteed by the Duplicate Lemma. 
Then, a straightforward induction shows
\begin{enumerate}
    \item  $|\langle \sigma \mid \sigma_{m^*}\rangle|^2 \ge 1 - 8^{m^*}\epsilon\ge 1-8k^3\epsilon$,
    \item  $\ket{\sigma_{m^*}}=\ket\phi^{\otimes k}$.
\end{enumerate}
That finishes the proof.
\end{proof}


We make a remark about \cref{thm:one_vs_many}. Note that some polynomial dependence on $k$ is unavoidable in this analysis for our procedure. Consider the following state:
\[
    \frac{1}{\sqrt {k+1}}\ket{\vec0}+\frac{1}{\sqrt {k+1}}\sum_{i=1}^{k}\ket{\vec{e_i}}. 
\]
To obtain a tight version of the above theorem with linear dependency on $k$ is an interesting problem.


\iffalse
Now we use our new slicing de Finetti theorem to prove the lemma used in the previous section.
\begin{lemma}
Let $\rho, \sigma $ be two states over some system $A_1\ldots A_2$. Suppose that 
SwapTest$(\rho, \sigma)$ passes w.p. $1-\epsilon$.
Then
\[
    \Exp_{\ket\sigma \in \rho} [\lambda_1(\ket\sigma)] \ge 1-O(\epsilon).
\]
\end{lemma}
\fi

\subsection{Many-versus-Many Slicing de Finetti}
In \cref{thm:one_vs_many}, we considered top Schmidt coefficient being large on a $1$ vs $k-1$ cut for pure state. By looking at the example we mentioned in the end of the previous subsection, it is natural to think that if the top Schmidt coefficient is large among a balanced cut, then we can obtain better trace distance. That is indeed the case. In fact, that top Schmidt coefficient is large for a balanced cut always implies the top Schmidt coefficient is large for a less balanced cut for a symmetric state. In this subsection, our goal is to formalize this intuition. 

\begin{theorem}[Many-versus-many Slicing de Finetti]\label{thm:many-many}
Let $\ket{\sigma}^{A_1\ldots A_k} \in \vee^k(\C^d)$. Suppose for some  $1\le\ell \le k/2$, the top Schmidt coefficient of $\ket\sigma$ over the $A_1\ldots A_\ell : A_{\ell+1}\ldots A_k$ cut is $\sqrt{1-\epsilon}$. Then there is  $\ket\phi\in\C_d$, such that
\[
    |\langle \sigma, \phi^{\otimes k}\rangle|^2 \ge 1 - O( (k/\ell)^3 \epsilon).
\]
\end{theorem}

We start by collecting a couple of useful facts. The first one says that if a symmetric state from $(\C^d)^{\otimes k}$ is close to a product state, then it is also close to a symmetric product state, i.e., $\ket{\phi^{\otimes k}}$ for some $\ket\phi \in \C ^d$.
\begin{lemma} \label{lem:prod-symmetric-state}
Given a symmetric state $\ket\sigma\in\Symm{k}{d}$ and a $k$-partite product state $\ket\psi\in (\C^d)^{\otimes k}$. Suppose $|\langle \psi \mid \sigma \rangle|^2 \ge 1-\epsilon$. Then there is $\ket\phi\in\C^d$ that satisfies
\[
    |\langle\sigma\mid \phi^{\otimes k}\rangle|^2 \ge 
    1-9\epsilon.
\]
\end{lemma}
\begin{proof}
We take advantage of $\ket\sigma$ being symmetric in a way similar to that of \cref{lem:dup}. As $\ket\sigma\in\Symm{k}{d}$, we have for any permutation $\pi\in \Sym_k$, $|\langle \sigma \mid \pi\psi\rangle|^2\ge1-\epsilon$. By \cref{fact:overlap-triangle-ineq},
\begin{align*}
    |\langle \psi \mid \pi\psi \rangle|^2 \ge 1-4\epsilon.
\end{align*}
Say $\ket\psi=\ket{\psi_1}\otimes\ket{\psi_2}\otimes \cdots \otimes\ket{\psi_k}$, then,
\begin{align}
    (1-4\epsilon)^{k!} 
    \le \prod_{\pi\in \Sym_k} |\langle \psi \mid \pi\psi \rangle|^2    
    %    \nonumber\\
    &=\left(\prod_{i\in[k]}\prod_{j \in [k]} |\langle \psi_i\mid \psi_j\rangle|^2\right)^{(k-1)!}
        \nonumber\\
    &\qquad\le \left(\Exp_{i\in[k]} \prod_{j\in[k]}
    |\langle \psi_i\mid \psi_j\rangle|^2\right)^{k!},
        \label{eq:prod-sym-ineq}
\end{align}
where the last step uses the AM-GM inequality. It follows from (\ref{eq:prod-sym-ineq}), there must exist $i\in[k]$ such that
\begin{align}
     1-4\epsilon\le \prod_{j\in [k]} |\langle \psi_i \mid \psi_j \rangle|^2 
         \quad \iff \quad
     1-4\epsilon \le |\langle  \psi_i ^{\otimes k} \mid \psi \rangle|^2.
     \nonumber
\end{align}
Apply \cref{fact:overlap-triangle-ineq} one more time, we obtain our lemma.
\end{proof}


The second fact  due to  Harrow and Montanaro~\cite[Appendix~B~Lemma~2]{HM13} and Soleimanifar and Wright~\cite{SW22matrix-product} establishes some criteria when a pure state is close to a product state.
\begin{lemma}
Given any quantum state $\ket\psi\in \cH_1 \otimes \cH_2 \otimes \cdots \otimes \cH_k$ for some arbitrary Hilbert space $\cH_1, \ldots, \cH_k$. Suppose 
\[
    \Exp_{\cS\subseteq [k]} [\trace{\psi_\cS^2}] \ge 1-\epsilon.
\]
Then for some product state $\ket\phi=\ket{\phi_1}\otimes\ket{\phi_2}\otimes\cdots\otimes\ket{\phi_k}$,
\[
    |\langle\psi \mid \phi\rangle|^2 \ge 1-3\epsilon.
\]
\end{lemma}
\iffalse
%  \psi 1-\delta close to product then
%  product test pass w.p. \le (1-\delta)^2 / 3  + 2 / 3 = 1 -2\delta/3 + \delta^2 / 3
%  let \eps = 2/3 \delta - \delta^/2 /3
%  SO, \psi 1 - \delta close to product ==> \psi 1-3\epsilon close to product
\fi
Combining the above two lemmas, we obtain
\begin{corollary}\label{cor:sym-prod}
Given any  state $\ket\sigma \in \Symm{k}{d}$. Suppose 
\[
    \Exp_{\cS\subseteq [k]} [\trace{\sigma_\cS^2}] \ge 1-\epsilon.
\]
Then for some  state $\ket\phi \in \C^d$,
\[
    |\langle\sigma \mid \phi^{\otimes k}\rangle|^2 \ge 1-27\epsilon.
\]
\end{corollary}

From the above discussion, to prove \cref{thm:many-many}, it suffices to bound $\Tr\psi_\cS^2$ for any subset $\cS$. The following ``cut lemma'' establishes such bounds. 
\begin{lemma}[Cut Lemma]\label{lem:cut}
Let $\ket{\sigma}\in \Symm{k}{d}$.
Suppose for some  $1\le\ell \le k/2$, the top Schmidt coefficient of $\ket\sigma$ over the $A_1\ldots A_\ell : A_{\ell+1}\ldots A_k$ cut is $\sqrt{1-\epsilon}$. Let $\cS\subseteq [k]$ be some arbitrary  subset.
Then, 
\[
    \trace \sigma_\cS ^2 \ge \begin{cases}
        1, & \quad |\cS|  = 0;
        \\
        1-6\epsilon, & \quad \min\{|\cS|, k-|\cS|\} \in \{1,2,\ldots, \ell-1\};
        \\
        1- O((|\cS|/\ell)^3\epsilon), &\quad \min\{|\cS| , k-|\cS|\}\in \{\ell,\ldots, k/2\}.
    \end{cases}
\]
\end{lemma}
\begin{proof}
For $\cS =\varnothing$, the statement is trivial as $\sigma$ is pure. Since $\ket\sigma\in\Symm{k}{d}$, without loss of generality, assume that 
$\cS=\{1, 2, \ldots, m\}$ for some $1 \le m \le k/2$. This is because 
$\trace \sigma_\cS^2 = \trace \sigma^2_{\bar\cS}$ when $\sigma$ is a pure state. Let $\ket\phi ^{A_1\ldots A_\ell} \ket\zeta ^{A_{\ell+1}\ldots A_k}$ be the 
top Schmidt component associated with the coefficient $\sqrt{1-\epsilon}$. 


\uline{Case 1}: $m < \ell$. Let $\cA = \{1, 2, \ldots, m\}, \cB = \{m+1, \ldots, \ell\}, \cC = \{\ell+1, \ldots, k-\ell+m \}, \cD = \{k-\ell+m+1, \ldots, k\}.$ Write down the Schmidt decomposition of $\ket\phi$ over the $\cA$ and $\cB$ cut, $\ket\zeta$ over the $\cC$ and $\cD$ cut,
\begin{equation}
    \ket\phi = \sum_i \sqrt{\lambda_i} \ket{\alpha_i}\ket{\beta_i}, 
    \qquad
    \ket\zeta = \sum_i \sqrt{\eta_i} \ket{\gamma_i} \ket{\delta_i}.
    \nonumber
\end{equation}
Since $\cB$ and $\cD$ has the same size, and that $\ket\sigma\in\Symm{k}{d}$, we have for the state $\ket\phi\ket\zeta$, if we switch the subsystem of $\cB$ and $\cD$, then the overlap with $\ket\sigma$ is still $1-\epsilon$. Therefore, by \cref{fact:overlap-triangle-ineq}, we have
\begin{align}
    (1-2\epsilon)^2 
    &\le \left|\left\langle
        \sum_{i,j} \sqrt{\lambda_i\eta_j}\bra{\alpha_i}^\cA\bra{\beta_i}^\cB
         \bra{\gamma_j}^\cC\bra{\delta_j}^\cD,
        %
        \sum_{i,j} \sqrt{\lambda_i\eta_j}\ket{\alpha_i}^\cA\ket{\beta_i}^\cD
        \ket{\gamma_j}^\cC\ket{\delta_j}^\cB
    \right\rangle\right|^2
    \nonumber \\
    & = \left(\sum_{i,j} \lambda_i\eta_j |\langle \beta_i \mid \delta_j \rangle|^2\right)^2
    \le  \lambda_1^2 \left(  \sum_{i,j}\eta_j |\langle \beta_i \mid \delta_j \rangle|^2 \right) ^2
    \nonumber\\
    &\le \lambda_1^2 \left(  \sum_{j}\eta_j \sum_i|\langle \beta_i \mid \delta_j \rangle|^2 \right)^2
    \le \lambda_1 ^2.
    \nonumber
\end{align}
Immediately,
\begin{align}
    \trace{\sigma_\cA ^2 } 
    &\ge (1-\epsilon)^2\trace [(\trace_{\cB\cC\cD}  (\phi\otimes\zeta) )^2]
    \nonumber \\
    &= (1-\epsilon)^2\trace \left[\left(\trace_{\cB\cC\cD} \left(
    \sum_i \lambda_i \alpha_i \otimes \beta_i\otimes\zeta
    \right) \right) ^2\right]
    \nonumber \\
    &\ge (1-\epsilon)^2 \lambda_1 ^2 \ge (1-\epsilon)^2 (1-2\epsilon)^2
    \nonumber\\
    &\ge 1-6\epsilon.\label{eq:many-many-case1-corr}
\end{align} 
The first step is true because $\sigma \succeq (1-\epsilon)\phi\otimes \zeta$, therefore  $\trace_{\cB\cC\cD}\sigma \succeq (1-\epsilon)\trace_{\cB\cC\cD}(\phi\otimes \zeta)$ as partial trace is completely positive. It then follows that $\trace\sigma_\cA ^2 \succeq (1-\epsilon)^2\trace(\trace_{\cB\cC\cD}(\phi\otimes\zeta)) ^2$.
\iffalse
\pnote{partial trace is completely positive}
\fi

\uline{Case 2}: $\ell < m \le k/2$. We are much like the situation of \cref{thm:one_vs_many}. Let $t = \lceil \log (m / \ell) \rceil$. For $i=1$ to $t$, we apply the Duplicate Lemma and obtain a state $\ket{\sigma_i}$, such that for $i=1,2,\ldots, t$
\begin{align}
    &\trace_{\{\ell\cdot 2^i+1,\ldots, k\}}\sigma_i = \phi ^{\otimes 2^i},
    \nonumber \\
    &\trace_{\{\ell\cdot 2^i+1,\ldots, k\}}\sigma_i^2 = 1,
    \label{eq:many-many-pure} \\
    &|\langle \sigma \mid \sigma_i \rangle|^2 \ge 1-8^i \epsilon.
    \label{eq:many-many-case2-corr}
\end{align}
By our choice of parameter, $2^{t-1}\ell < m \le 2^{t}\ell$. If $m=2^t \ell$, then $(\sigma_t)_\cS = \trace_{\ell\cdot2^t+1,\ldots,k}\sigma_t$ is pure by (\ref{eq:many-many-pure}). Then  
\begin{align*}
    \sqrt {\trace\sigma_\cS^2} =  \sqrt{\trace \sigma_\cS ^2 \cdot \trace (\sigma_t)_\cS^2} &\ge \trace(\sigma_\cS\cdot (\sigma_t)_\cS) = F(\sigma_\cS, (\sigma_t)_\cS) \\
    &\ge F(\sigma, \sigma_t)  = |\langle \sigma \mid \sigma_t\rangle|^2 \ge (1-8^{\log(m/\ell)})\epsilon,
\end{align*}
where the first step and third step are true because $(\sigma_t)_\cS$ is pure; the second step uses (\ref{eq:matrix-cauchy-schwarz}); the fourth step is by \cref{fact:data-processing}, the data processing inequality for fidelity; then fifth step is again by purity of the states; and the final step uses (\ref{eq:many-many-case2-corr}). It follows that
\[
    \trace \sigma_\cS^2 \ge 1 - O((m/\ell)^3\epsilon).
\]
If $m<2^t\ell$, then we can apply Case 1. 
Let $\cA= \{1,2,\ldots,2^t\ell\}, \cB = \{2^t\ell+1, \ldots, k\}$. Then in view of (\ref{eq:many-many-case2-corr}), the top Schmidt coefficient of $\ket\sigma$ among the $\cA : \cB$ cut is at least $\sqrt{1-8^t\epsilon}$ by \cref{claim:top-schmidt-coeff}. Thus by (\ref{eq:many-many-case1-corr}),
\[
    \trace \sigma_\cS^2 \ge 1-6\cdot 8^t\epsilon\ge 1-O((m/\ell)^3\epsilon).
\qedhere
\]
% one can do better, but asymptotically nothing improves
\end{proof}

Now \cref{thm:many-many} follows from \cref{cor:sym-prod} and \cref{lem:cut}.




\subsection{Proof of \cref{thm:de-Finnetti-type-1}}
Now we record a version of the slicing de Finetti theorem for the mixture of symmetric states. A natural generalization of the top Schmidt coefficient among some $A:B$ cut for a state $\sigma$ being large is that $\trace{\sigma_A^2} $ being large.
In particular,
\begin{lemma}\label{lem:trace-reduced-den-squared}
Let $\sigma\in\C^n\otimes\C^m$ be some density operator, and $A,B$ are the systems with respect to the space  $\C^n$ and $\C^m$, respectively. Suppose
\begin{equation*}
    \trace \sigma_A^2 \ge 1-\epsilon.
\end{equation*}
Let $\mu$ be some distribution on pure states induced by $\sigma$, then
\begin{align*}
    \Exp_{\rho\sim\mu} \lambda_1(\rho) \ge 1-\epsilon.
\end{align*}
\end{lemma}
\begin{proof}
Let $m=|\supp\mu|$ be a finite number, this is without loss of generality.  Let $\rho_1, \rho_2,\ldots, \rho_m$ be the pure states in $\supp\mu$. Further, write the Schmidt decomposition for each $\rho_i$
\begin{align*}
    \ket{\rho_i} = \sum_{j} \sqrt{\lambda_{ij}}\ket{\phi_{ij}}^{A}\ket{\sigma_{ij}}^{B},\qquad \lambda_{i1} \ge \lambda_{i2}\ge\cdots.
\end{align*}
Then
\begin{align*}
    \sigma_A = \sum_i \mu(\rho_i) \sum_j \lambda_{ij} \ket{\phi_{ij}}\bra{\phi_{ij}}.
\end{align*}
Thus,
\begin{align}
    \trace{\sigma_{A}^2} 
    &= \sum_i \mu(\rho_i)^2 \sum_j \lambda_{ij}^2 +    \sum_{i\not=i'}\mu(\rho_i)\mu(\rho_{i'})\sum_{j,j'}\lambda_{ij}\lambda_{i'j'}|\langle \phi_{ij} \mid \phi_{i'j'}\rangle|^2
    \nonumber\\
    &\le \sum_i \mu(\rho_i)^2 \sum_j \lambda_{ij}^2 + \sum_{i\not=i'}\mu(\rho_i)\mu(\rho_{i'})\lambda_{i1}\sum_{j,j'}\lambda_{i'j'}|\langle \phi_{ij} \mid \phi_{i'j'}\rangle|^2
    \nonumber\\
    &\le \sum_i \mu(\rho_i)^2 \sum_j \lambda_{ij}^2 + \sum_{i\not=i'}\mu(\rho_i)\mu(\rho_{i'})\lambda_{i1}\sum_{j'}\lambda_{i'j'}
    \nonumber\\
    &\le \sum_i \mu(\rho_i)^2  \sum_j \lambda_{ij}^2 + \sum_{i\not=i'}\mu(\rho_i)\mu(\rho_{i'}) \lambda_{i1}
    \nonumber\\
    &=\sum_i \mu(\rho_i)^2\sum_j \lambda_{ij}^2 + \sum_i \mu(\rho_i)(1-\mu(\rho_i))\lambda_{i1}
    \nonumber\\
    &\le \sum_i \mu(\rho_i)^2 \lambda_{i1} + \sum_i \mu(\rho_i)(1-\mu(\rho_i))\lambda_{i1}
    \nonumber\\
    &= \sum_i \mu(\rho_i)\lambda_{i1},
    \nonumber
\end{align}
where the third step holds because for fixed $i,i',j'$,  $\sum_{j} |\langle\phi_{ij}\mid \phi_{i'j'}\rangle|^2\le 1.$
\end{proof}


\begin{theorem}\label{thm:conv-symm-de-Finetti}
Given density operator $\sigma^{A_1\ldots A_k}$ that describes states from $\conv{\Symm{k}{d}}$. 
For any $1\le \ell\le k/2$ and $\cA=\{A_1,A_2,\cdots, A_\ell\}$, there is some distribution $\mu$ on $\ket\phi\in\C^d$,
\begin{align}
    \left\|
        \sigma - \int \ket\phi\bra\phi^{\otimes k} d\mu
    \right\|_1 \le O\left(\sqrt{(k/\ell)^3(1-\trace{\sigma_{\cA}^2}) }\right).
\end{align}

\end{theorem}
\begin{proof}
Let $\mu$ be the distribution on pure symmetric states induced by $\sigma$. Let $\trace\sigma_{\cA}^2=1-\epsilon$. The theorem follows immediately by combining \cref{fact:trace-innerproduct}, \cref{lem:trace-reduced-den-squared}, \cref{thm:many-many}, and triangle inequality.
\end{proof}
